% Shared exercise chunk: l3-probabilistic-classification
\begin{exercisechunk}
\item \textbf{Classification from probabilistic prediction.}
\label[exercise]{ex:probabilistic-classification}
Let $(X,Y)$ have an arbitrary distribution on $\cX\times\{0,1\}$, where
$\cX$ is finite. For every $x$ with $\P(X=x)>0$, define
\[
\eta(x)=\P(Y=1\mid X=x).
\]
At the remaining inputs, $\eta(x)$ may be defined arbitrarily. For a decision
$a\in\{0,1\}$ and an $x$ with $\P(X=x)>0$, define the conditional risk
\[
r_x(a)=\P(a\ne Y\mid X=x).
\]
For a classifier $h:\cX\to\{0,1\}$, define its risk by
\[
R_{01}(h)=\P(h(X)\ne Y).
\]
A classifier is \emph{Bayes-optimal} if it minimizes $R_{01}$ over all
classifiers $h:\cX\to\{0,1\}$. The minimum value
\[
R_{01}^*=\min_{h:\cX\to\{0,1\}}R_{01}(h)
\]
is the \emph{Bayes risk}.
\begin{enumerate}[(a)]
\item Compute $r_x(0)$ and $r_x(1)$. Show that
\[
h^*(x)=\I\{\eta(x)\ge1/2\}
\]
is Bayes-optimal and that the Bayes risk is
\[
R_{01}^*=\mathbb E[\min\{\eta(X),1-\eta(X)\}].
\]
\item Prove the exact excess-risk identity
\[
R_{01}(h)-R_{01}^*
=\mathbb E\left[
|2\eta(X)-1|\I\{h(X)\ne h^*(X)\}
\right].
\]
\item Let $q:\cX\to(0,1)$ be a probabilistic
predictor, and define
\[
L_{\log}(q)
=\mathbb E[-Y\log q(X)-(1-Y)\log(1-q(X))].
\]
Let $L_{\log}^*=L_{\log}(\eta)$, with the usual conventions at
$\eta(x)\in\{0,1\}$, and let $h_q(x)=\I\{q(x)\ge1/2\}$. Prove
\[
R_{01}(h_q)-R_{01}^*
\le \sqrt{2\bigl(L_{\log}(q)-L_{\log}^*\bigr)}.
\]
\item \emph{Optional extension.} Construct a family of probabilistic
predictors whose induced
classifiers $h_q$ are Bayes-optimal but whose excess log loss tends to
infinity.
\end{enumerate}

\emph{Hint for (c):} You may use Pinsker's inequality for Bernoulli
distributions:
\[
|p-q|\le
\sqrt{\frac12\left(
p\log\frac pq+(1-p)\log\frac{1-p}{1-q}
\right)}.
\]

\paragraph{What this exercise is meant to teach.}
A probabilistic predictor and a classifier answer different questions. The
excess classification risk weights a wrong decision by twice the distance of
$\eta(x)$ from the decision boundary $1/2$. Small excess log loss controls
excess classification risk, but Bayes-optimal classification does not imply
accurate probabilities or small log loss.
\end{exercisechunk}
